\section*{الفصل \ref{ch:g6:exploring-our-environment} --- استكشاف بيئتنا}

\begin{solution}{exo:g6:exploring-our-environment:1}
\omterm{def:g6:exploring-our-environment:components}{الكائنات الحية} (الطحلب، وحشرات القمل الخشبي)؛ \omterm{def:g6:exploring-our-environment:components}{والمادة المعدنية}
(حجر الجدار، وبركة الماء)؛ وآثار \omterm{def:g6:exploring-our-environment:components}{الكائنات الحية} وبقاياها (\omterm{def:g1:plants-around-us:parts}{الأوراق}
الميتة، \omterm{def:g1:animals-around-us:covering}{وصدفة} الحلزون).
\end{solution}

\begin{solution}{exo:g6:exploring-our-environment:2}
\omterm{def:g6:exploring-our-environment:conditions}{درجة الحرارة}، بمحرار (بالدرجات المئوية)؛ \omterm{def:g6:exploring-our-environment:conditions}{ورطوبة} الهواء والتربة؛
\omterm{def:g6:exploring-our-environment:conditions}{والضوء}، من الشمس الكاملة إلى الظل العميق.
\end{solution}

\begin{solution}{exo:g6:exploring-our-environment:3}
الحي: السرخسة، وأم أربعة وأربعين. والمعدني: بركة الماء، وحجر
الجدار. والآثار: \omterm{def:g1:animals-around-us:covering}{صدفة} الحلزون، وفضلات \omterm{prop:g3:sorting-living-things:groups}{الطيور}.
\end{solution}

\begin{solution}{exo:g6:exploring-our-environment:4}
ستّ عشرة درجة (\qty{38}{\degreeCelsius} في مقابل
\qty{22}{\degreeCelsius}) وشمس كاملة في مقابل ظل عميق (وجفاف في
مقابل ندى أيضًا). وعشبة الجفاف والوهج توافق القدم الجنوبية؛
والطحلب \omterm{prop:g4:plants-without-seeds:damp}{والسرخسيات} وحشرات القمل الخشبي توافق الركن البارد الندي.
\end{solution}

\begin{solution}{exo:g6:exploring-our-environment:5}
كساؤها يدع الماء يفلت، ففي شمس جافة يفقد جسمها ماءه فقدًا قاتلًا
سريعًا. ولهذا توجد حشرات القمل الخشبي في \omterm{def:g2:where-animals-live:habitat}{موائل} دقيقة ندية مظلمة ---
تحت الحجارة وتحت الخشب الميت.
\end{solution}

\begin{solution}{exo:g6:exploring-our-environment:6}
مثلًا: بندقة مقضومة --- فأر الحرش؛ وشبكة عنكبوت --- عنكبوتها؛
وآثار أقدام في الطين --- البلشون (والفضلات \omterm{def:g1:animals-around-us:covering}{والريش} والأصداف تصلح
الصلاح نفسه).
\end{solution}

\begin{solution}{exo:g6:exploring-our-environment:7}
صندوق مغلق نصفه ورق ندي ونصفه ورق جاف، وحشرات قمل خشبي تُطلق على
خط الوسط؛ وبعد دقائق تجلس كلها تقريبًا على الجانب الندي (وكذلك
الظل في مقابل \omterm{def:g6:exploring-our-environment:conditions}{الضوء}). وهذا يدل على أن كل \omterm{def:g1:animals-around-us:animal}{حيوان} يستقر حيث تناسب
\omterm{def:g6:exploring-our-environment:conditions}{الشروط} جسمه --- فخريطة الجماعة ترسم نفسها بلا سَوق.
\end{solution}

\begin{solution}{exo:g6:exploring-our-environment:8}
لأن القياسات المأخوذة بالطريقة نفسها هي وحدها القابلة للمقارنة:
فالفرق يجب أن يأتي من المواقع لا من القياس. وهي النسخة الميدانية
من الاختبار المنصف --- شيء واحد يتغير (الموقع)، وكل ما عداه يبقى
واحدًا.
\end{solution}

\begin{solution}{exo:g6:exploring-our-environment:9}
أعلى الحجر: مضاء، جاف، دافئ --- أشنة (والذبابة المتشمسة زائرة).
ووجهه السفلي: مظلم، ندي، بارد --- حشرات قمل خشبي (أو أم أربعة
وأربعين). خمسة سنتيمترات من الفارق، ومجموعتان متوافقتان من
السكان.
\end{solution}

\begin{solution}{exo:g6:exploring-our-environment:10}
الخطوة الأولى: علّم مترًا مربعًا محددًا. والثانية: قِس --- قراءة
المحرار مكتوبة، وندى التربة محسوسًا مسجّلًا، ودرجة \omterm{def:g6:exploring-our-environment:conditions}{الضوء} مسجّلة.
والثالثة: فتّش كما ينبغي، بحجارة تُرفع وتُعاد، وبأصناف المكوّنات
الثلاثة كلها معدَّدة. والرابعة: سجّل في الموقع --- أعدادًا لا
``بعضًا''. والخامسة: لم تكن المقارنة ممكنة أصلًا بموقع واحد؛ فامسح
موقعًا ثانيًا بالطريقة نفسها.
\end{solution}

\begin{solution}{exo:g6:exploring-our-environment:11}
التفسير: اللبلاب يزدهر في مدى \omterm{def:g6:exploring-our-environment:conditions}{شروط} الجدار الشمالي (الظل والندى) لا
في مدى \omterm{def:g6:exploring-our-environment:conditions}{شروط} الجنوبي (الوهج والجفاف). والاختبار: قِس قدمي الجدارين
--- الحرارة والندى \omterm{def:g6:exploring-our-environment:conditions}{والضوء} --- وراجع تفضيلات اللبلاب المعروفة؛ أو
جد اللبلاب على جدران أخرى وسجّل أي اتجاه تشترك فيه جدرانه.
\end{solution}

\begin{solution}{exo:g6:exploring-our-environment:12}
أزل فرق الجدار: أجرِ الاختيار في صندوق مستدير (فالجدار واحد في كل
مكان)، مع تبديل النصفين الندي والجاف بين المرتين --- الندى شرقًا في
المرة الأولى، وغربًا في الثانية. فإن تبعت \omterm{prop:g3:sorting-living-things:groups}{الحشرات} الندى حيثما
وُضع، خرجت الجدران من الحساب؛ ولم يتغير بين المرتين إلا الشرط
المختبَر، وبقي كل ما عداه واحدًا.
\end{solution}

\begin{solution}{pb:g6:exploring-our-environment:1}
\textbf{1.} السياج النباتي العالي: فهو يظلّل الخندق الشمالي
(ويحميه)، أما الجنوبي فيواجه الشمس طول النهار.

\textbf{2.} \omterm{def:g6:exploring-our-environment:conditions}{درجة الحرارة}: \qty{19}{\degreeCelsius} في مقابل
\qty{31}{\degreeCelsius}. \omterm{def:g6:exploring-our-environment:conditions}{والرطوبة}: تربة الشمالي ندية (تلتصق)،
وتربة الجنوبي جافة (غبار). \omterm{def:g6:exploring-our-environment:conditions}{والضوء} يكمل الثلاثي: ظل في مقابل شمس
كاملة.

\textbf{3.} \omterm{def:g6:exploring-our-environment:conditions}{الشروط} تتغير عبر النهار: ففي السادسة مساءً يكون أي
موقع أبرد وضوؤه أخفض، فيخلط الفرق المقيس الموقعَ بالساعة. والساعة
نفسها تُبقي المقارنة منصفة.

\textbf{4.} ليتساوى جهد التفتيش: فمساحة أكبر أو زمن أطول في موقع
يجد فيه سكانًا أكثر --- وهو اختبار غير منصف لأي الخندقين أغنى.

\textbf{5.} القائمتان لا تسمّيان إلا كائنات حية. والناقص:
المكوّنات المعدنية (التربة، والحجارة، والماء) والآثار --- وهي مما
يجب أن يسجّله الجرد الصحيح أيضًا.

\textbf{6.} حشرات القمل الخشبي: الخندق الشمالي --- فهي تحتاج إلى
الندى والظل وإلا جفّت. والعلجوم: الشمالي --- فجلده رطب ويحتاج إلى
مأوى ندي. والسحلية: الجنوبي --- فهي تدفّئ جسمها في الشمس، وتصطاد
حيث تستطيع التشمس.

\textbf{7.} إما أن الحلزونات تعيش هناك مختفية (ولعلها لا تنشط إلا
في ساعات الليل الندية)، وإما أن \omterm{def:g1:animals-around-us:covering}{الصدفة} بقية من مكان آخر أو من موسم
آخر --- فالأثر يثبت الحضور \emph{في وقت ما}، لا الإقامة الآن.

\textbf{8.} مظلم، ندي، بارد. والقاعدة هي قاعدة
\cref{prop:g6:exploring-our-environment:decide}: كل بقعة يسكنها
\omterm{def:g5:biodiversity:species}{النوع} الذي توافق شروطها مداه.

\textbf{9.} في الشمالي طحلب وعلاجيم لأن ظله يُبقيه باردًا نديًّا،
وهو ما تقتضيه أجسامها؛ وفي الجنوبي سحالٍ وأعشاب ذات \omterm{def:g1:plants-around-us:parts}{أوراق} شمعية
لأن شمسه الكاملة تجعله حارًّا جافًّا، وهو ما يناسب صائدًا متشمسًا
ونباتًا مقاومًا للجفاف.

\textbf{10.} من غير السياج ينال الخندق الشمالي شمسًا كاملة: فيصير
أحرّ وأجفّ --- \omterm{def:g6:exploring-our-environment:conditions}{بشروط} تنزع نحو \omterm{def:g6:exploring-our-environment:conditions}{شروط} الخندق الجنوبي اليوم. فتوقّع
أن يتضاءل محبّو الندى (الطحلب، \omterm{prop:g4:plants-without-seeds:damp}{والسرخسيات}، وحشرات القمل الخشبي،
والعلجوم) وأن ينتقل إليه محبّو الشمس (الأعشاب الجافة، والجنادب،
وربما السحلية).

\textbf{11.} لأن المسح الثاني يجب أن يختلف عن الأول في شيء واحد
فحسب --- غياب السياج. فمواقع أو أدوات أو موسم جديدة تضيف فروقًا
خاصة بها، فلا يعود التغير منسوبًا إلى السياج: وهي قاعدة الاختبار
المنصف على مقياس منظر طبيعي.

\textbf{12.} كل حشرة قمل خشبي تهيم وتستدير حتى يرتاح جسمها، فتنتهي
الجماعة مجتمعة حيث الندى --- فمواضع \omterm{def:g1:animals-around-us:animal}{الحيوانات} تسجّل \omterm{def:g6:exploring-our-environment:conditions}{الشروط}، بيقين
قراءات المحرار نفسه. ولا أحد يسوقها؛ بل راحة كل فرد هي التي ترسم
الخريطة، وهذا بالضبط ما يجعل السكان قابلين للقراءة قياساتٍ.
\end{solution}
